\begin{figure}[h]

\usepackage{neurips_2026}
\usepackage[utf8]{inputenc}
\usepackage[T1]{fontenc}
\usepackage{hyperref}
\usepackage{url}
\usepackage{booktabs}
\usepackage{graphicx}
\usepackage{amsmath}
\usepackage{xcolor}
\usepackage{microtype}

\graphicspath{{figures/}{auto/}}

% Placeholder macro: any number not yet measured renders red with a dagger.
\newcommand{\ph}[1]{\textcolor{red}{\underline{#1}$^{\dagger}$}}

% AUTO-GENERATED numbers. Every quantity quoted in this paper resolves through
% autopilot/p8_make_assets.py, which reads only verified JSON artifacts. No
% number below is typed by hand, so prose, tables and figures cannot disagree.
% AUTO-GENERATED by autopilot/p8_make_assets.py - do not edit by hand
% Every number in the manuscript resolves through these macros.
\newcommand{\Ntest}{3000}
\newcommand{\Npos}{1466}
\newcommand{\Nneg}{1534}
\newcommand{\Nboot}{10{,}000}
\newcommand{\Nprobes}{23}
\newcommand{\AUCRandomEpFifty}{0.8641}
\newcommand{\CIloRandomEpFifty}{0.8510}
\newcommand{\CIhiRandomEpFifty}{0.8769}
\newcommand{\AUCOracleEpFifty}{0.8740}
\newcommand{\CIloOracleEpFifty}{0.8615}
\newcommand{\CIhiOracleEpFifty}{0.8862}
\newcommand{\AUCEnvelopeEpFifty}{0.8761}
\newcommand{\CIloEnvelopeEpFifty}{0.8635}
\newcommand{\CIhiEnvelopeEpFifty}{0.8879}
\newcommand{\DOracleRandomEpFifty}{+0.0099}
\newcommand{\DOracleRandomEpFiftyCI}{[+0.0051,\,+0.0149]}
\newcommand{\DOracleRandomEpFiftyP}{$<$0.0001}
\newcommand{\DOracleRandomEpFiftyQ}{0.0001}
\newcommand{\DEnvelopeRandomEpFifty}{+0.0120}
\newcommand{\DEnvelopeRandomEpFiftyCI}{[+0.0068,\,+0.0173]}
\newcommand{\DEnvelopeRandomEpFiftyP}{$<$0.0001}
\newcommand{\DEnvelopeRandomEpFiftyQ}{$<$0.0001}
\newcommand{\DOracleEnvelopeEpFifty}{-0.0020}
\newcommand{\DOracleEnvelopeEpFiftyCI}{[-0.0069,\,+0.0029]}
\newcommand{\DOracleEnvelopeEpFiftyP}{0.4114}
\newcommand{\DOracleEnvelopeEpFiftyQ}{0.4114}
\newcommand{\AUCRandomEpSeventyFive}{0.8723}
\newcommand{\CIloRandomEpSeventyFive}{0.8593}
\newcommand{\CIhiRandomEpSeventyFive}{0.8849}
\newcommand{\AUCOracleEpSeventyFive}{0.8836}
\newcommand{\CIloOracleEpSeventyFive}{0.8714}
\newcommand{\CIhiOracleEpSeventyFive}{0.8955}
\newcommand{\AUCEnvelopeEpSeventyFive}{0.8803}
\newcommand{\CIloEnvelopeEpSeventyFive}{0.8677}
\newcommand{\CIhiEnvelopeEpSeventyFive}{0.8922}
\newcommand{\DOracleRandomEpSeventyFive}{+0.0113}
\newcommand{\DOracleRandomEpSeventyFiveCI}{[+0.0063,\,+0.0164]}
\newcommand{\DOracleRandomEpSeventyFiveP}{$<$0.0001}
\newcommand{\DOracleRandomEpSeventyFiveQ}{$<$0.0001}
\newcommand{\DEnvelopeRandomEpSeventyFive}{+0.0080}
\newcommand{\DEnvelopeRandomEpSeventyFiveCI}{[+0.0028,\,+0.0132]}
\newcommand{\DEnvelopeRandomEpSeventyFiveP}{0.0025}
\newcommand{\DEnvelopeRandomEpSeventyFiveQ}{0.0045}
\newcommand{\DOracleEnvelopeEpSeventyFive}{+0.0033}
\newcommand{\DOracleEnvelopeEpSeventyFiveCI}{[-0.0013,\,+0.0079]}
\newcommand{\DOracleEnvelopeEpSeventyFiveP}{0.1491}
\newcommand{\DOracleEnvelopeEpSeventyFiveQ}{0.1677}
\newcommand{\AUCRandomEpHundred}{0.8746}
\newcommand{\CIloRandomEpHundred}{0.8618}
\newcommand{\CIhiRandomEpHundred}{0.8869}
\newcommand{\AUCOracleEpHundred}{0.8855}
\newcommand{\CIloOracleEpHundred}{0.8735}
\newcommand{\CIhiOracleEpHundred}{0.8971}
\newcommand{\AUCEnvelopeEpHundred}{0.8807}
\newcommand{\CIloEnvelopeEpHundred}{0.8683}
\newcommand{\CIhiEnvelopeEpHundred}{0.8926}
\newcommand{\DOracleRandomEpHundred}{+0.0109}
\newcommand{\DOracleRandomEpHundredCI}{[+0.0057,\,+0.0160]}
\newcommand{\DOracleRandomEpHundredP}{$<$0.0001}
\newcommand{\DOracleRandomEpHundredQ}{$<$0.0001}
\newcommand{\DEnvelopeRandomEpHundred}{+0.0062}
\newcommand{\DEnvelopeRandomEpHundredCI}{[+0.0010,\,+0.0113]}
\newcommand{\DEnvelopeRandomEpHundredP}{0.0199}
\newcommand{\DEnvelopeRandomEpHundredQ}{0.0299}
\newcommand{\DOracleEnvelopeEpHundred}{+0.0047}
\newcommand{\DOracleEnvelopeEpHundredCI}{[+0.0004,\,+0.0091]}
\newcommand{\DOracleEnvelopeEpHundredP}{0.0294}
\newcommand{\DOracleEnvelopeEpHundredQ}{0.0378}
\newcommand{\AUCAnatomyTwoEpFifty}{0.8654}
\newcommand{\DAnatomyTwoRandomEpFifty}{+0.0013}
\newcommand{\DAnatomyTwoRandomEpFiftyCI}{[-0.0054,\,+0.0082]}
\newcommand{\DAnatomyTwoRandomEpFiftyP}{0.7101}
\newcommand{\AUCCoverEpFifty}{0.8643}
\newcommand{\DCoverRandomEpFifty}{+0.0002}
\newcommand{\DCoverRandomEpFiftyCI}{[-0.0049,\,+0.0053]}
\newcommand{\DCoverRandomEpFiftyP}{0.9440}
\newcommand{\AUCAncestor}{0.8487}
\newcommand{\NBlack}{431}
\newcommand{\NWhite}{2318}
\newcommand{\NAsian}{251}
\newcommand{\NprobesRace}{19}
\newcommand{\WorstRaceCount}{19}
\newcommand{\BlackAUCOracle}{0.8472}
\newcommand{\BlackCIOracle}{[0.8100,\,0.8804]}
\newcommand{\WhiteAUCOracle}{0.8853}
\newcommand{\RaceGapOracle}{0.0382}
\newcommand{\GapRho}{+0.511}
\newcommand{\GapRhoP}{0.0255}
\newcommand{\SexGapRho}{-0.714}
\newcommand{\SexGapRhoP}{0.0006}
\newcommand{\NprobesSub}{19}
\newcommand{\SubGenderRho}{-0.714}
\newcommand{\SubGenderRhoP}{0.0006}
\newcommand{\SubGenderQ}{0.0042}
\newcommand{\SubGenderBranchRho}{-0.857}
\newcommand{\SubGenderBranchP}{0.0137}
\newcommand{\SubGenderGapMin}{0.0272}
\newcommand{\SubGenderGapMax}{0.0396}
\newcommand{\SubGenderWorst}{female}
\newcommand{\SubGenderWorstN}{19}
\newcommand{\SubRaceRho}{+0.511}
\newcommand{\SubRaceRhoP}{0.0255}
\newcommand{\SubRaceQ}{0.0595}
\newcommand{\SubRaceBranchRho}{+0.429}
\newcommand{\SubRaceBranchP}{0.3374}
\newcommand{\SubRaceGapMin}{0.0475}
\newcommand{\SubRaceGapMax}{0.0935}
\newcommand{\SubRaceWorst}{black}
\newcommand{\SubRaceWorstN}{19}
\newcommand{\SubEthnicityRho}{+0.351}
\newcommand{\SubEthnicityRhoP}{0.1408}
\newcommand{\SubEthnicityQ}{0.1971}
\newcommand{\SubEthnicityBranchRho}{+0.464}
\newcommand{\SubEthnicityBranchP}{0.2939}
\newcommand{\SubEthnicityGapMin}{0.0223}
\newcommand{\SubEthnicityGapMax}{0.0661}
\newcommand{\SubEthnicityWorst}{hispanic}
\newcommand{\SubEthnicityWorstN}{19}
\newcommand{\SubLanguageRho}{+0.275}
\newcommand{\SubLanguageRhoP}{0.2537}
\newcommand{\SubLanguageQ}{0.2537}
\newcommand{\SubLanguageBranchRho}{+0.786}
\newcommand{\SubLanguageBranchP}{0.0362}
\newcommand{\SubLanguageGapMin}{0.0518}
\newcommand{\SubLanguageGapMax}{0.0805}
\newcommand{\SubLanguageWorst}{other languages}
\newcommand{\SubLanguageWorstN}{15}
\newcommand{\SubMaritalRho}{-0.444}
\newcommand{\SubMaritalRhoP}{0.0570}
\newcommand{\SubMaritalQ}{0.0997}
\newcommand{\SubMaritalBranchRho}{-0.714}
\newcommand{\SubMaritalBranchP}{0.0713}
\newcommand{\SubMaritalGapMin}{0.0364}
\newcommand{\SubMaritalGapMax}{0.0654}
\newcommand{\SubMaritalWorst}{single}
\newcommand{\SubMaritalWorstN}{19}
\newcommand{\SubAgeRho}{+0.304}
\newcommand{\SubAgeRhoP}{0.2065}
\newcommand{\SubAgeQ}{0.2409}
\newcommand{\SubAgeBranchRho}{+0.357}
\newcommand{\SubAgeBranchP}{0.4316}
\newcommand{\SubAgeGapMin}{0.0487}
\newcommand{\SubAgeGapMax}{0.0704}
\newcommand{\SubAgeWorst}{$<$60}
\newcommand{\SubAgeWorstN}{18}
\newcommand{\SubSeverityRho}{-0.521}
\newcommand{\SubSeverityRhoP}{0.0222}
\newcommand{\SubSeverityQ}{0.0595}
\newcommand{\SubSeverityBranchRho}{-0.857}
\newcommand{\SubSeverityBranchP}{0.0137}
\newcommand{\SubSeverityGapMin}{0.1286}
\newcommand{\SubSeverityGapMax}{0.1400}
\newcommand{\SubSeverityWorst}{mild ($-6$,$-2$]}
\newcommand{\SubSeverityWorstN}{19}
\newcommand{\NbranchesSub}{7}
\newcommand{\SeverityGapSpread}{0.0114}
\newcommand{\SubAUCMin}{0.8483}
\newcommand{\SubAUCMax}{0.8855}
\newcommand{\SevRandomMild}{0.8164}
\newcommand{\SevRandomModerate}{0.9030}
\newcommand{\SevRandomSevere}{0.9525}
\newcommand{\SevOracleMild}{0.8301}
\newcommand{\SevOracleModerate}{0.9131}
\newcommand{\SevOracleSevere}{0.9588}
\newcommand{\SevEnvelopeMild}{0.8232}
\newcommand{\SevEnvelopeModerate}{0.9103}
\newcommand{\SevEnvelopeSevere}{0.9559}
\newcommand{\SevDeltaMild}{+0.0137}
\newcommand{\SevDeltaModerate}{+0.0102}
\newcommand{\SevDeltaSevere}{+0.0063}
\newcommand{\SevGapRandom}{0.1361}
\newcommand{\SevGapOracle}{0.1286}
\newcommand{\RaceRandomWhite}{0.8741}
\newcommand{\RaceRandomBlack}{0.8325}
\newcommand{\RaceRandomAsian}{0.9033}
\newcommand{\RaceOracleWhite}{0.8853}
\newcommand{\RaceOracleBlack}{0.8472}
\newcommand{\RaceOracleAsian}{0.9189}
\newcommand{\BlackGainOracle}{+0.0147}
\newcommand{\FTOracleMeanpool}{0.8947}
\newcommand{\FTRandomMeanpool}{0.8868}
\newcommand{\FTDelta}{+0.0079}
\newcommand{\FTOracleBest}{0.8947}
\newcommand{\FTRandomBest}{0.8878}
\newcommand{\FTDeltaBest}{+0.0069}


\title{Location Beats Shape: A Controlled Comparison of Anatomy-Guided\\
Masking for Joint-Embedding Predictive Pretraining\\
on Retinal OCT}

\author{}

\begin{document}

\maketitle

\begin{abstract}
Masked predictive pretraining asks a model to reconstruct the representation of
hidden image regions. Which regions are hidden is a free design choice, and in
medical imaging there is an intuitive answer: hide the tissue that carries the
diagnosis. We test that intuition under tightly controlled conditions. Six
masking policies are continued from a single shared I-JEPA checkpoint on 3D
retinal OCT, differing \emph{only} in how predictor targets are placed, and
evaluated with a frozen linear-probe protocol for glaucoma classification
(FairVision, $N{=}\Ntest$). Guidance does help: restricting targets to retinal
tissue improves AUC by $\DEnvelopeRandomEpFifty$ over unguided masking at the
matched epoch (95\% CI $\DEnvelopeRandomEpFiftyCI$), and a band located by a
first-order intensity statistic improves it by $\DOracleRandomEpHundred$ at
epoch 100 ($\DOracleRandomEpHundredCI$, $p\,\DOracleRandomEpHundredP$). All six
contrasts against the null, at all three epochs, exclude zero. But
anatomical \emph{precision} does not help. A policy that places $97\%$ of masked
patches on tissue, using a trained retinal segmentation model to shape the
targets, is \emph{indistinguishable from unguided masking} at the matched epoch
($\DAnatomyTwoRandomEpFifty$, CI $\DAnatomyTwoRandomEpFiftyCI$), as is a
coverage-constrained policy ($\DCoverRandomEpFifty$) --- while the two policies
that merely \emph{place} rectangles on tissue gain an order of magnitude more.
Carried past epoch 50 the coverage-constrained policy does not merely stall, it
falls \emph{behind} unguided masking ($\DCoverRandomEpSeventyFive$ at epoch 75,
CI $\DCoverRandomEpSeventyFiveCI$), so over-constraining where the model may look
is worse than not guiding it at all. Direct measurement of mask geometry across
600 slices shows the fraction of anatomy hidden is uncorrelated with downstream
AUC. The best policy consults no segmentation model at all. A subgroup audit over
\NprobesSub{} probes finds the same worst-served group every time; gains reach
every group and every disease stage, mild disease most of all
($\SevDeltaMild$), but the difficulty ordering is untouched: the severe-to-mild
gap stays within $\SeverityGapSpread$ of constant across every policy. Because
each policy was pretrained once, we report these as single-run observations
rather than an established ranking.
\end{abstract}

\section{Introduction}

Self-supervised learning has become the default route to representation quality
in medical imaging, where expert annotation is expensive and pathology is
rare~\citep{azizi2021big,zhou2023retfound,qiu2024visionfm}. Joint-embedding
predictive architectures such as I-JEPA~\citep{assran2023ijepa} learn by
predicting the \emph{representation} of masked target blocks from a visible
context block, avoiding pixel reconstruction and its bias toward high-frequency
detail~\citep{he2022mae,xie2022simmim}.

I-JEPA samples its targets geometrically: rectangles placed uniformly at random.
In natural images this is defensible, since informative content is spread
broadly across the frame. In a retinal OCT B-scan it is not obviously so. A
large fraction of the image is dark vitreous and background; the diagnostic
signal for glaucoma is concentrated in a thin band of retinal tissue. This
motivates an intuitive hypothesis, one that several recent methods
adopt~\citep{li2022semmae,li2024anatomask,wang2025maskwhatmatters,shi2022adios}:

\begin{quote}
\emph{Direct the predictive objective at clinically meaningful anatomy rather
than at arbitrary image regions.}
\end{quote}

We test this hypothesis directly and at several strengths. Rather than proposing
one anatomy-aware method and reporting that it beats a baseline, we construct a
\emph{ladder} of masking policies of increasing anatomical specificity --- from
unguided rectangles, through rectangles constrained to lie on tissue, to ragged
targets whose shape follows the segmented retina --- and hold everything else
fixed. All arms continue from one shared pretrained checkpoint, use identical
optimiser, schedule, and effective batch size, and are evaluated under one
frozen-probe protocol.

The result contradicts the intuition in an informative way. Guidance helps, but
the benefit does not increase with anatomical precision --- it disappears. The
two policies whose targets track the segmented retina most faithfully are
statistically indistinguishable from unguided masking at the matched epoch, while
two policies that merely place ordinary rectangles on tissue gain an order of
magnitude more. The best performer uses no segmentation model at all, locating a
band by a per-column intensity centroid.

\paragraph{Contributions.}
\begin{itemize}
\item A controlled six-arm study isolating masking policy in medical SSL, with a
      shared ancestor checkpoint and a single evaluation protocol (21 frozen
      probes).
\item Evidence that guidance helps but anatomical precision does not, including the
      specific control --- rectangles restricted to retinal tissue --- that
      separates ``avoid background'' from ``target anatomy''.
\item Direct measurement of mask geometry for every policy on 600 slices,
      showing that fraction-of-anatomy-hidden does not order downstream AUC, and
      an explicit accounting of the axes on which the anatomy arms are
      confounded.
\item A subgroup audit and a $10^{-5}$ reproduction of the headline result from
      public weights on different hardware and precision.
\end{itemize}

\section{Related work}

\paragraph{Masked image modelling and JEPAs.}
Masked autoencoders reconstruct pixels~\citep{he2022mae,xie2022simmim};
BEiT~\citep{bao2022beit} and iBOT~\citep{zhou2022ibot} predict discrete tokens;
data2vec~\citep{baevski2022data2vec} and I-JEPA~\citep{assran2023ijepa} predict
latent representations. Video and 3D extensions follow the same
template~\citep{bardes2024vjepa,chen2023mim3d,taleb20203dssl}. Recent analyses
question how much of the benefit derives from the masking distribution
itself~\citep{balestriero2025lejepa,he2025dseqjepa}.

\paragraph{Informed masking.}
A substantial line of work argues that \emph{what} is masked matters.
ADIOS~\citep{shi2022adios} learns an adversarial masking model;
SemMAE~\citep{li2022semmae} uses learned semantic parts;
AttMask~\citep{kakogeorgiou2022attmask} masks by attention;
HardPatches~\citep{wang2023hardpatches} and AutoMAE~\citep{chen2023automae}
target difficulty. In the medical domain, AnatoMask~\citep{li2024anatomask},
AnatomyMAE~\citep{ceballosarroyo2026anatomymae},
MSMAE~\citep{mao2023msmae}, and VasoMIM~\citep{huang2026vasomim} explicitly
steer masking toward anatomy, and \citet{wang2025maskwhatmatters} argue directly
for masking clinically relevant regions. Our study is a controlled test of that
family's central premise on one task, and finds it does not hold in the direction
usually assumed.

\paragraph{Evidence that informed masking is not uniformly better.}
The picture is less one-sided than the paragraph above suggests. The original
masked-modelling papers already report that random sampling is hard to beat:
MAE's ablation puts random-75 at $73.5$ linear against block-75 at
$63.9$~\citep{he2022mae}, and SimMIM finds random beats its best block
variant~\citep{xie2022simmim}. Among informed schemes, AttMask sits within $0.2$
points of random~\citep{kakogeorgiou2022attmask}, AutoMAE ties MAE under full
fine-tuning~\citep{chen2023automae}, HPM gains only when randomness is retained
and \emph{loses} when restricted to the hardest patches~\citep{wang2023hpm}, and
attention-guided MAE underperforms MAE in the low-label
regime~\citep{sick2025attg}. Most directly, SemMAE's own part-quality ablation
shows that \emph{imperfect} semantic parts add nothing over no parts at
all~\citep{li2022semmae}, and \citet{guo2025mlomae} reproduce SemMAE below random
MAE under a common protocol, with the deficit largest under frozen evaluation ---
also our protocol. \citet{chen2023automae} name the mechanism a ``patch selection
dilemma'': a slight foreground bias helps, but raising it too far starves the
encoder of the evidence it must predict from. Conversely, the closest positive
precedent for our best policy uses no segmentation model either:
\citet{lee2025hufgm} select foreground from normalised intensity alone and gain
across five 3D medical benchmarks. Appendix~\ref{app:counter} tabulates these.

We therefore do \emph{not} claim to overturn the informed-masking literature.
That random masking is a strong baseline, and semantic guidance wins only
sometimes, is known and contested. What we add is a controlled measurement of
where a \emph{trained medical segmentation model} falls on that spectrum when
everything else is held fixed, in a setting --- 3D retinal OCT, joint-embedding
prediction, frozen probe --- where we could find no prior controlled comparison.

\paragraph{Retinal foundation models and OCT.}
RETFound~\citep{zhou2023retfound}, VisionFM~\citep{qiu2024visionfm},
URFound~\citep{yu2024urfound}, and OCTCube~\citep{liu2026octcube} pretrain on
large retinal corpora; OCT-specific masked modelling is explored by
\citet{pissas2024octmim}. MIRAGE~\citep{morano2025mirage} provides multimodal
OCT segmentation and supplies the anatomical guidance used by several of our
arms. Deep glaucoma detection from OCT is
established~\citep{maetschke2019glaucoma,ran2019glaucoma,noury2022glaucoma}.

\paragraph{Fairness in medical imaging.}
Performance disparities across demographic groups are widely
documented~\citep{seyyedkalantari2021underdiagnosis,larrazabal2020gender,gichoya2022race}.
FairVision~\citep{luo2024fairvision} and Harvard-GF~\citep{luo2024harvardgf}
provide OCT data with demographic attributes;
FairMedFM~\citep{jin2024fairmedfm} and \citet{shi2025equitable} benchmark
fairness for medical foundation models. We report a subgroup analysis as a
secondary finding.

\section{Method: a ladder of masking policies}
\label{sec:policies}

\subsection{Background}

I-JEPA partitions the token grid of an image into a context block and $M$ target
blocks. A context encoder $f_\theta$ embeds the visible tokens; a predictor
$g_\phi$ maps that embedding plus positional queries to predicted
representations of the target tokens; an exponential-moving-average target
encoder $f_{\bar\theta}$ supplies the regression targets. With a
$16{\times}16$ token grid the design choice we vary is which of the 256 cells
become targets.

\subsection{The six policies}

All policies produce $M{=}4$ target blocks and one context block under identical
scale and aspect-ratio ranges. They differ only in placement and shape.

\begin{description}
\item[\textsc{random} (null).] Stock I-JEPA multiblock. Four rectangles placed
  uniformly at random anywhere in the grid. No image content is consulted.

\item[\textsc{envelope} (random-within-retina).] Identical rectangles ---
  same shape, size and count as \textsc{random} --- but rejection-sampled onto
  the retinal envelope predicted by MIRAGE~\citep{morano2025mirage}. Only
  \emph{location} changes. This is the control that separates ``stop wasting
  targets on background'' from ``target anatomy specifically''.

\item[\textsc{intensity} (segmentation-free).] A band whose vertical position is
  located per slice by a per-column intensity-weighted row centroid, smoothed
  across columns. It adapts to the retina's position and curvature using a
  first-order intensity statistic, and consults \textbf{no segmentation model}.
  It uses no ground-truth annotation of any kind. (Our released code and
  checkpoints label this policy \texttt{oracle}, a historical name from when it
  was intended as an upper bound. We rename it here because it is neither an
  oracle nor an upper bound: it reads only the input image.)

\item[\textsc{anatomy-v1} / \textsc{anatomy-v2}.] Ragged connected components
  grown on the MIRAGE score map, so target \emph{shape} follows tissue rather
  than being rectangular. v2 additionally bridges diagonal adjacencies.

\item[\textsc{cover} $f$.] Targets chosen greedily to cover the anatomy subject
  to a hard floor leaving a fraction $f$ of tissue visible in the context. We
  pretrain $f{=}0.21$.
\end{description}

Figure~\ref{fig:policies} shows all six on the same B-scans.

\begin{figure}[t]
\centering
\includegraphics[width=0.74\linewidth]{fig1_policies_compact.png}
\caption{The six masking policies on one identical B-scan, drawn with the
production samplers. Each predictor target block has its own colour; the green
contour is the anatomy reference. Anatomical specificity increases along the
ladder: \textsc{envelope} and \textsc{cover} place rectangles on tissue, while
the anatomy arms grow targets that trace the retina almost exactly. Downstream
AUC does not increase with that specificity (Table~\ref{tab:main}).
Figure~\ref{fig:policies5} repeats this over five B-scans spanning the volume.}
\label{fig:policies}
\end{figure}

\subsection{Hypotheses}

The ladder yields three testable statements:

\begin{description}
\item[H1] \emph{Avoiding background helps.} \textsc{envelope} $>$ \textsc{random}.
\item[H2] \emph{Anatomical precision helps further.} anatomy arms $>$
  \textsc{envelope}.
\item[H3] \emph{Any benefit is due to anatomical targeting, not to incidental
  changes in mask area, target count, or context size.}
\end{description}

\section{Experimental setup}
\label{sec:setup}

\paragraph{Data.} FairVision~\citep{luo2024fairvision} glaucoma OCT volumes.
Each volume is a 3D scan resampled to 100 B-scans of $256{\times}256$ after
transform. The label is a binary glaucoma diagnosis. Splits are fixed and
patient-disjoint as released; downstream evaluation uses a held-out test split of
$N{=}3000$ volumes (1466 positive / 1534 negative), seed 42, fixed loader order.
No test volume is seen during pretraining, probe fitting, or model selection: the
probe head is selected on a separate validation split. Labels are byte-identical
across arms, so any two arms are paired-comparable on the same cases. The dataset
carries self-reported demographic attributes, which we use only for the subgroup
analysis in Section~\ref{sec:subgroup} and never as model inputs.

\paragraph{Pretraining.} Patch-level I-JEPA, ViT-Base, $256{\times}256$ crops,
patch size 16, predictor depth 6, EMA target encoder. \textbf{Every arm resumes
from the same epoch-25 checkpoint} and continues with identical optimiser,
learning-rate and weight-decay schedules, and effective batch size (512 for all
arms). Guidance is ramped linearly from epoch 25 to 30 and held at full strength
thereafter. The \emph{only} difference between arms is the mask sampler.

Arms differ in how far they were carried, and we do not pretend otherwise.
\textsc{random}, \textsc{intensity} and \textsc{envelope} reach epoch 100.
\textsc{anatomy-v2} has valid probes to epoch 50 (later checkpoints follow a
target-precision splice and are excluded, Appendix~\ref{app:excluded}).
\textsc{anatomy-v1} has only epoch 30. \textsc{cover} $f{=}0.21$ was
deliberately halted at epoch 73. Epoch 50 is therefore the only epoch at which
five policies are simultaneously available, and it carries every cross-policy
comparison that involves the anatomy or coverage arms.

\paragraph{Evaluation.} The encoder is frozen. Per-slice features are
mean-pooled over the volume, followed by LayerNorm and a linear head. The head is
selected by validation AUC and reported on the held-out test split. We report
paired bootstrap confidence intervals over test volumes (\Nboot{} resamples,
stratified by class, the same resampled index set applied to every arm) and
DeLong tests for correlated ROC curves~\citep{delong1988roc}, both of which are
valid here because all arms are evaluated on the same cases in the same order.
Each test volume is a distinct subject, so case-level resampling is already
subject-level and no cluster correction is
required~\citep{obuchowski1997clusteredroc,rutter2000clusterbootstrap}. We follow
\citet{maierhein2024metricsreloaded} in reporting the operating-point metrics
alongside AUC rather than AUC alone (Appendix~\ref{app:allprobes}).

\paragraph{Numerical precision, and how we handle it.} The probe harness
defaults to mixed precision. The \textsc{random}, \textsc{intensity} and
\textsc{envelope} long-horizon probes were fitted under that default (fp16);
the \textsc{anatomy-v1}, \textsc{anatomy-v2}, \textsc{cover} and ancestor probes
were fitted with autocast explicitly disabled (fp32). We therefore do
\emph{not} claim a single protocol across every probe in the study. Instead we
partition the comparisons: all headline contrasts in
Table~\ref{tab:main} are drawn from arms sharing both epoch \emph{and}
precision, and any contrast that would cross the precision boundary is marked
as such and excluded from headline claims. Section~\ref{sec:precision} reports a
direct measurement of how large the precision effect actually is.

\section{Results}

\emph{What these numbers can and cannot establish.} Each policy was pretrained
exactly once, so policy is perfectly confounded with one stochastic optimisation
path after the fork. The intervals below quantify sampling error over test
subjects for a \emph{fixed} pair of score vectors, not seed-to-seed variance of
pretraining. Read every comparison here as a controlled single-run observation,
not an established ranking over retrainings (Section~\ref{sec:limits}).

\subsection{H1 holds, H2 fails}

\begin{table}[t]
\centering
\caption{Frozen MeanPool probe, FairVision test ($N{=}\Ntest$, \Npos{} positive).
All arms continue from one epoch-25 ancestor (AUC $\AUCAncestor$). $\Delta$ is
versus \textsc{random} at the \emph{matched} epoch. The upper block is
precision-matched (fp16 throughout) and carries every headline claim. The lower
block is fp32. Its epoch-50 $\Delta$ is taken against an fp32 null re-probed for
that comparison and is matched on epoch and precision. The \textsc{cover}
epoch-75 $\Delta$ is marked $\ddagger$: no fp32 null exists at that epoch yet, so
it is matched on epoch only. The measured fp16/fp32 gap is at most
$2\times10^{-4}$ (Appendix~\ref{app:fp32}), some forty times below that
$\Delta$.}
\label{tab:main}
\footnotesize
\setlength{\tabcolsep}{4pt}
\begin{tabular}{llcccccc}
\toprule
& & \multicolumn{2}{c}{epoch 50} & \multicolumn{2}{c}{epoch 75} & \multicolumn{2}{c}{epoch 100} \\
\cmidrule(lr){3-4}\cmidrule(lr){5-6}\cmidrule(lr){7-8}
policy & guidance signal & AUC & $\Delta$ & AUC & $\Delta$ & AUC & $\Delta$ \\
\midrule
\textsc{random}   & none (null)                      & \AUCRandomEpFifty  & ---                              & \AUCRandomEpSeventyFive  & ---                              & \AUCRandomEpHundred  & ---  \\
\textsc{envelope} & segmenter (location only)        & \AUCEnvelopeEpFifty & $\mathbf{\DEnvelopeRandomEpFifty}$ & \AUCEnvelopeEpSeventyFive & $\mathbf{\DEnvelopeRandomEpSeventyFive}$ & \AUCEnvelopeEpHundred & $\mathbf{\DEnvelopeRandomEpHundred}$ \\
\textsc{intensity}   & intensity centroid, no segmenter & \AUCOracleEpFifty   & $\mathbf{\DOracleRandomEpFifty}$   & \AUCOracleEpSeventyFive   & $\mathbf{\DOracleRandomEpSeventyFive}$   & \textbf{\AUCOracleEpHundred} & $\mathbf{\DOracleRandomEpHundred}$ \\
\midrule
\textsc{anatomy-v2} & segmenter (target shape) & \AUCAnatomyTwoEpFifty & \DAnatomyTwoRandomEpFifty & --- & --- & --- & --- \\
\textsc{cover} $f{=}.21$ & segmenter (coverage) & \AUCCoverEpFifty & \DCoverRandomEpFifty & \AUCCoverEpSeventyFive & $\mathbf{\DCoverRandomEpSeventyFive}^{\ddagger}$ & \ph{pending} & \ph{pending} \\
\bottomrule
\end{tabular}
\end{table}


Table~\ref{tab:main} gives the main comparison, Table~\ref{tab:contrasts} in
Appendix~\ref{app:contrasts} the full paired inference for all nine
precision-matched contrasts, and Figure~\ref{fig:traj} the trajectories. All six
contrasts against the null survive Benjamini--Hochberg correction over that
family of nine ($q \leq \DEnvelopeRandomEpHundredQ$ throughout).
\textbf{H1 holds}: constraining the same rectangles to lie on retinal tissue
(\textsc{envelope}) beats unguided masking by $\DEnvelopeRandomEpFifty$ at epoch
50 (CI $\DEnvelopeRandomEpFiftyCI$) and by $\DEnvelopeRandomEpHundred$ at epoch
100 (CI $\DEnvelopeRandomEpHundredCI$). Avoiding background is worth something,
at every epoch we measured, with an interval that excludes zero.

\begin{figure}[t]
\centering
\includegraphics[width=0.95\linewidth]{fig_trajectories_ci.png}
\caption{\textbf{(a)} Frozen-probe test AUC against pretraining epoch. All arms
fork from the same epoch-25 checkpoint (open circle) and differ only in masking
policy. Solid lines are the precision-matched fp16 family; dashed lines are fp32
arms and must not be compared vertically against the solid lines.
\textbf{(b)} The actual inference: paired differences against the unguided null
on identical test cases, with \Nboot-resample bootstrap intervals. All six
intervals exclude zero. We plot paired differences rather than per-arm error
bars because marginal intervals carry between-case variance that cancels in a
paired comparison, and would understate the evidence.}
\label{fig:traj}
\end{figure}

\textbf{H2 fails.} The arms that use the segmentation model most aggressively ---
shaping targets to trace the retina, or maximising anatomy coverage under a
floor --- do not improve on the null. At epoch 50, the only epoch at which all
five arms have a probe, \textsc{anatomy-v2} gains $+0.0013$ over
\textsc{random} and \textsc{cover} gains $\DCoverRandomEpFifty$, against
$\DEnvelopeRandomEpFifty$ for \textsc{envelope} and $\DOracleRandomEpFifty$ for
\textsc{intensity}. The policies that merely \emph{place} rectangles on tissue gain
roughly an order of magnitude more than the policies that shape targets to match
it. Increasing anatomical fidelity past ``rectangles on tissue'' does not add
benefit.

We are careful not to overstate this. At the matched epoch 50 neither
\textsc{anatomy-v2} (CI $\DAnatomyTwoRandomEpFiftyCI$,
$p=\DAnatomyTwoRandomEpFiftyP$) nor \textsc{cover} (CI
$\DCoverRandomEpFiftyCI$, $p=\DCoverRandomEpFiftyP$) is \emph{worse} than the
null; both are indistinguishable from it, and both are measured against an fp32
null re-probed specifically for this comparison, so they are matched on epoch
\emph{and} precision. What the epoch-50 data rule out is not that these policies
are harmful, but the expected \emph{monotone improvement with anatomical
precision}: their intervals are centred near zero and comfortably contain it,
while the two cruder policies produce intervals that exclude it.

\paragraph{Carried further, coverage-constrained masking stalls and falls
behind.} \textsc{cover} is the only segmenter-driven policy we continued past
epoch 50, and it produces the study's one \emph{negative} result. Its trajectory
is $\AUCCoverEpTwentySeven$, $\AUCCoverEpThirty$, $\AUCCoverEpThirtyFour$,
$\AUCCoverEpFifty$, $\AUCCoverEpSeventyThree$, $\AUCCoverEpSeventyFive$: it
climbs normally to epoch 50 and then stops. From epoch 50 to 75 it moves
$\DCoverSelfFiftyToSeventyFive$ ($p=\DCoverSelfFiftyToSeventyFiveP$), while the
null gains $\DRandomSelfFiftyToSeventyFive$. By epoch 75 it sits \emph{below} the
null by $\DCoverRandomEpSeventyFive$ (CI $\DCoverRandomEpSeventyFiveCI$,
$p=\DCoverRandomEpSeventyFiveP$), below \textsc{envelope} by
$\DCoverEnvelopeEpSeventyFive$ and below \textsc{intensity} by
$\DCoverOracleEpSeventyFive$. Forcing a fixed fraction of tissue to stay visible
does not merely fail to help; it appears to cap what the encoder can still learn
once the easy signal is exhausted. Two caveats: these probes are fp32 against
fp16 comparators, though the measured precision effect is at most
$2\times10^{-4}$ (Appendix~\ref{app:fp32}), some forty times below this gap; and
only $f{=}0.21$ was pretrained, so this concerns that floor, not a dose-response.

The strongest policy is the one that uses no segmentation model.
\textsc{intensity} gives $\DOracleRandomEpHundred$ over the null at epoch 100
(95\% CI $\DOracleRandomEpHundredCI$, $p\,\DOracleRandomEpHundredP$,
$q=\DOracleRandomEpHundredQ$), with a margin that is stable across training
($\DOracleRandomEpFifty$, $\DOracleRandomEpSeventyFive$, $\DOracleRandomEpHundred$). It exceeds
the segmenter-guided \textsc{envelope} arm at epoch 100 by
$\DOracleEnvelopeEpHundred$ (CI $\DOracleEnvelopeEpHundredCI$,
$p=\DOracleEnvelopeEpHundredP$), though the two are indistinguishable at epochs 50
and 75. Notably, \textsc{envelope}'s margin over the null \emph{decays} with
training ($\DEnvelopeRandomEpFifty \to \DEnvelopeRandomEpSeventyFive \to
\DEnvelopeRandomEpHundred$) while \textsc{intensity}'s does not.

\subsection{H3: it is not anatomical targeting}
\label{sec:results-h3}

\begin{table}[t]
\centering
\caption{Measured mask geometry, 600 held-out slices, production samplers.
``purity'' is the fraction of masked cells lying on tissue; ``loss slots'' is the
number of predictor targets contributing to the loss. All values measured. AUC is
quoted at the \emph{matched} epoch 50, the only epoch at which all five policies
have a probe, so that the geometry-to-AUC association is not read across mixed
epochs. All five AUCs are at epoch 50.}
\label{tab:geom}
\small
\begin{tabular}{lcccccc}
\toprule
policy & anatomy hidden & purity & mask ratio & context kept & loss slots & AUC @ep50 \\
\midrule
\textsc{random}   & 52.2\% & 31.6\% & 43.7\% & 42.1\% & 157.7 & \AUCRandomEpFifty \\
\textsc{intensity}   & 62.2\% & 41.1\% & 40.0\% & 45.6\% & 158.4 & \AUCOracleEpFifty \\
\textsc{envelope} & 76.9\% & 43.5\% & 46.4\% & 40.5\% & 159.9 & \textbf{\AUCEnvelopeEpFifty} \\
\textsc{cover}    & 74.1\% & 45.3\% & 43.3\% & 43.5\% & 160.0 & \AUCCoverEpFifty \\
\textsc{anatomy-v2} & 80.3\% & \textbf{97.3\%} & 21.4\% & 67.9\% & 64.0 & \AUCAnatomyTwoEpFifty \\
\bottomrule
\end{tabular}
\end{table}

To test H3 we ran each production sampler on 600 held-out slices and measured
what it actually does (Table~\ref{tab:geom}, Figure~\ref{fig:paradox}). Three
observations follow.

\begin{figure}[t]
\centering
\includegraphics[width=0.46\linewidth]{fig_precision_paradox.png}
\caption{Downstream AUC at matched epoch 50 against mask purity, the fraction of
masked patches lying on retinal tissue. The policy whose masks are $97\%$
on-tissue does not separate from the null, while policies at $41$--$44\%$ purity
gain the most. The ordering does not support the intuition that more accurate
anatomical masking yields better representations. With four to five arms this is
descriptive, not inferential; the same result arranged as a specificity ladder is
in Appendix~\ref{app:fp32}.}
\label{fig:paradox}
\end{figure}
